\documentclass[11pt,a4paper]{article}
\usepackage[T1]{fontenc}
\usepackage[utf8]{inputenc}
\usepackage{amsmath,amssymb,xcolor}
\usepackage[margin=25mm]{geometry}
\usepackage[colorlinks=true,urlcolor=blue]{hyperref}
\definecolor{nexorablue}{HTML}{2459B5}
\setlength{\parindent}{0pt}
\setlength{\parskip}{7pt}
\setlength{\emergencystretch}{2em}
\renewcommand{\familydefault}{\sfdefault}
\begin{document}

{\small\bfseries\color{nexorablue} NEXORA CPGE}\par
{\small MPSI / NEX-ENS-03 / Corrigé}\par
{\LARGE\bfseries Quand deux ensembles disjoints ont une image commune}\par
\vspace{0.25em}\hrule\vspace{0.6em}

\section*{Images directes}
Sur $[-1,0]$, $|x-1|$ parcourt $[1,2]$ en décroissant ; sur $[2,3]$, il parcourt le même segment en croissant. Les deux images valent donc $[1,4]$.

\section*{Antécédents}
L’inégalité sur la valeur absolue donne $[-1,0]\cup[2,3]$. Les deux segments de départ sont disjoints, leur intersection a une image vide, tandis que l’intersection des images vaut $[1,4]$.


\end{document}
